\section{Redes convolucionales}
\label{sec:marco_vision}

Clasificar una imagen consiste en asignarla a una de $C$ categorías: la entrada es la imagen y la salida, el
vector de probabilidades $p_c$ de la Ecuación~\ref{eq:marco_softmax}. Para eso hay que convertir una matriz
de píxeles, de dimensión muy alta, en un vector de características que la cabeza de clasificación pueda
separar por clase. Las \glspl{cnn}, consolidadas a fines de la década de 1990, lo hacen componiendo filtros
locales~\cite{lecun1998gradient}, y fueron la arquitectura dominante en clasificación de imágenes hasta la
llegada de los transformers.

\subsection{Convolución y arquitectura de una red convolucional}
\label{sec:marco_cnn}

La operación de convolución desliza un filtro de pequeño tamaño (\emph{kernel}) $K$ de $k \times k$ pesos
sobre la entrada $I$ y calcula en cada posición el producto punto entre el filtro y la ventana que cubre,
\begin{equation}
    S(i, j) = \sum_{m=0}^{k-1}\sum_{n=0}^{k-1} I(i+m,\, j+n)\, K(m, n).
    \label{eq:marco_convolucion}
\end{equation}
El resultado $S$ es un mapa de características, que vale alto donde la ventana se parece al patrón que el
filtro codifica. La conectividad local supone que la información relevante se agrupa en regiones pequeñas, y
el uso de un mismo filtro en toda la imagen, que un patrón puede aparecer en cualquier parte
(Figura~\ref{fig:marco_convolucion}). Una capa aplica muchos filtros a la vez, cada uno con sus propios pesos,
y apila sus mapas como canales de la salida~\cite{lecun1998gradient,goodfellow2016deep}.

% Figura: la operación de convolución.
% Se usa en sec:marco_cnn, después del párrafo que define la convolución.
% Las tres matrices van separadas y las etiquetas en dos líneas, para que no se toquen (tutor, p36).

\begin{figure}[H]
    \centering
    \begin{tikzpicture}[x=1mm, y=1mm]

        % Entrada 6x6 con la ventana de 3x3 resaltada.
        \fill[fignaranja!30] (5, 10) rectangle (20, 25);
        \draw[figtrazo!50] (0, 0) grid[step=5] (30, 30);
        \draw[figtrazo, thick] (0, 0) rectangle (30, 30);
        \draw[fignaranja!70!black, very thick] (5, 10) rectangle (20, 25);
        \node[font=\footnotesize, text=figtrazo, align=center, anchor=north] at (15, -2) {entrada\\$6\times 6$};

        % Filtro 3x3.
        \fill[fignaranja!30] (50, 7.5) rectangle (65, 22.5);
        \draw[fignaranja!70!black, shift={(50, 7.5)}] (0, 0) grid[step=5] (15, 15);
        \draw[fignaranja!70!black, very thick] (50, 7.5) rectangle (65, 22.5);
        \node[font=\footnotesize, text=figtrazo, align=center, anchor=north] at (57.5, -2) {filtro $3\times 3$\\(pesos compartidos)};
        \node[font=\large, text=figtrazo] at (40, 15) {$\ast$};

        % Salida 4x4 con la celda resultante.
        \fill[figverde!35] (95, 10) rectangle (100, 15);
        \draw[figtrazo!50] (90, 0) grid[step=5] (110, 20);
        \draw[figtrazo, thick] (90, 0) rectangle (110, 20);
        \draw[figverde!70!black, very thick] (95, 10) rectangle (100, 15);
        \node[font=\footnotesize, text=figtrazo, align=center, anchor=north] at (100, -2) {mapa de características\\$4\times 4$};
        \node[font=\large, text=figtrazo] at (77.5, 15) {$=$};

        % De la ventana a la celda.
        \draw[figflecha, fignaranja!70!black, bend left=25] (20, 25) to
            node[figetiqueta, above] {producto punto} (95, 15);

    \end{tikzpicture}
    \caption{Convolución: un mismo filtro recorre la entrada y produce un mapa de características.}
    \label{fig:marco_convolucion}
\end{figure}


Una \gls{cnn} alterna capas de convolución, activaciones no lineales y capas de reducción (\emph{pooling}),
que resumen cada ventana de un mapa, por ejemplo con su máximo, y reducen la resolución. Cada reducción hace que
una unidad de la capa siguiente dependa de una región más grande de la imagen, su campo receptivo. Por eso las
primeras capas responden a bordes, las siguientes a esquinas y contornos y las últimas a partes de
objetos~\cite{goodfellow2016deep}. Al final, un promedio global de cada canal convierte los mapas en un vector
de características, y la cabeza de clasificación de la Ecuación~\ref{eq:marco_capa} lo convierte en los
logits~\cite{he2015residual}.

Apilar más capas agranda el campo receptivo, pero a partir de cierta profundidad aparece la degradación:
agregar capas aumenta incluso el error de entrenamiento, de modo que el problema no es el sobreajuste
(Sección~\ref{sec:marco_regularizacion}). ResNet la resuelve con conexiones residuales: cada bloque aprende una
corrección $F$ que se suma a su entrada,
\begin{equation}
    \mathbf{h}^{(l+1)} = \mathbf{h}^{(l)} + F\bigl(\mathbf{h}^{(l)}\bigr),
    \label{eq:marco_residual}
\end{equation}
de modo que, si la identidad es lo mejor para un bloque, le basta con llevar $F$ a cero, y se pueden entrenar
redes de más de cien capas~\cite{he2015residual}. DenseNet lleva la idea más lejos: cada capa recibe la concatenación de
los mapas de todas las capas anteriores del bloque, $\mathbf{h}^{(l)} = H_l\bigl([\mathbf{h}^{(0)}, \dots,
\mathbf{h}^{(l-1)}]\bigr)$, con lo que reutiliza los rasgos ya calculados y alcanza una capacidad comparable con
menos parámetros~\cite{huang2018densely}. ConvNeXt rediseña una ResNet con decisiones tomadas de los
transformers: convoluciones por canal de $7 \times 7$, normalización por capa, un bloque que expande y vuelve a
reducir los canales y menos activaciones. Con esos cambios, una red puramente convolucional iguala o supera en exactitud
al Swin Transformer de la Sección~\ref{sec:marco_swin}~\cite{liu2022convnet}.

Una convolución solo relaciona píxeles cercanos, y vincular dos regiones lejanas de la imagen exige muchas
capas. La atención, en cambio, relaciona cualquier par de posiciones en una sola capa.
